% Options for packages loaded elsewhere
\PassOptionsToPackage{unicode}{hyperref}
\PassOptionsToPackage{hyphens}{url}
\documentclass[
]{article}
\usepackage{xcolor}
\usepackage{amsmath,amssymb}
\setcounter{secnumdepth}{-\maxdimen} % remove section numbering
\usepackage{iftex}
\ifPDFTeX
  \usepackage[T1]{fontenc}
  \usepackage[utf8]{inputenc}
  \usepackage{textcomp} % provide euro and other symbols
\else % if luatex or xetex
  \usepackage{unicode-math} % this also loads fontspec
  \defaultfontfeatures{Scale=MatchLowercase}
  \defaultfontfeatures[\rmfamily]{Ligatures=TeX,Scale=1}
\fi
\usepackage{lmodern}
\ifPDFTeX\else
  % xetex/luatex font selection
\fi
% Use upquote if available, for straight quotes in verbatim environments
\IfFileExists{upquote.sty}{\usepackage{upquote}}{}
\IfFileExists{microtype.sty}{% use microtype if available
  \usepackage[]{microtype}
  \UseMicrotypeSet[protrusion]{basicmath} % disable protrusion for tt fonts
}{}
\makeatletter
\@ifundefined{KOMAClassName}{% if non-KOMA class
  \IfFileExists{parskip.sty}{%
    \usepackage{parskip}
  }{% else
    \setlength{\parindent}{0pt}
    \setlength{\parskip}{6pt plus 2pt minus 1pt}}
}{% if KOMA class
  \KOMAoptions{parskip=half}}
\makeatother
\usepackage{color}
\usepackage{fancyvrb}
\newcommand{\VerbBar}{|}
\newcommand{\VERB}{\Verb[commandchars=\\\{\}]}
\DefineVerbatimEnvironment{Highlighting}{Verbatim}{commandchars=\\\{\}}
% Add ',fontsize=\small' for more characters per line
\newenvironment{Shaded}{}{}
\newcommand{\AlertTok}[1]{\textcolor[rgb]{1.00,0.00,0.00}{\textbf{#1}}}
\newcommand{\AnnotationTok}[1]{\textcolor[rgb]{0.38,0.63,0.69}{\textbf{\textit{#1}}}}
\newcommand{\AttributeTok}[1]{\textcolor[rgb]{0.49,0.56,0.16}{#1}}
\newcommand{\BaseNTok}[1]{\textcolor[rgb]{0.25,0.63,0.44}{#1}}
\newcommand{\BuiltInTok}[1]{\textcolor[rgb]{0.00,0.50,0.00}{#1}}
\newcommand{\CharTok}[1]{\textcolor[rgb]{0.25,0.44,0.63}{#1}}
\newcommand{\CommentTok}[1]{\textcolor[rgb]{0.38,0.63,0.69}{\textit{#1}}}
\newcommand{\CommentVarTok}[1]{\textcolor[rgb]{0.38,0.63,0.69}{\textbf{\textit{#1}}}}
\newcommand{\ConstantTok}[1]{\textcolor[rgb]{0.53,0.00,0.00}{#1}}
\newcommand{\ControlFlowTok}[1]{\textcolor[rgb]{0.00,0.44,0.13}{\textbf{#1}}}
\newcommand{\DataTypeTok}[1]{\textcolor[rgb]{0.56,0.13,0.00}{#1}}
\newcommand{\DecValTok}[1]{\textcolor[rgb]{0.25,0.63,0.44}{#1}}
\newcommand{\DocumentationTok}[1]{\textcolor[rgb]{0.73,0.13,0.13}{\textit{#1}}}
\newcommand{\ErrorTok}[1]{\textcolor[rgb]{1.00,0.00,0.00}{\textbf{#1}}}
\newcommand{\ExtensionTok}[1]{#1}
\newcommand{\FloatTok}[1]{\textcolor[rgb]{0.25,0.63,0.44}{#1}}
\newcommand{\FunctionTok}[1]{\textcolor[rgb]{0.02,0.16,0.49}{#1}}
\newcommand{\ImportTok}[1]{\textcolor[rgb]{0.00,0.50,0.00}{\textbf{#1}}}
\newcommand{\InformationTok}[1]{\textcolor[rgb]{0.38,0.63,0.69}{\textbf{\textit{#1}}}}
\newcommand{\KeywordTok}[1]{\textcolor[rgb]{0.00,0.44,0.13}{\textbf{#1}}}
\newcommand{\NormalTok}[1]{#1}
\newcommand{\OperatorTok}[1]{\textcolor[rgb]{0.40,0.40,0.40}{#1}}
\newcommand{\OtherTok}[1]{\textcolor[rgb]{0.00,0.44,0.13}{#1}}
\newcommand{\PreprocessorTok}[1]{\textcolor[rgb]{0.74,0.48,0.00}{#1}}
\newcommand{\RegionMarkerTok}[1]{#1}
\newcommand{\SpecialCharTok}[1]{\textcolor[rgb]{0.25,0.44,0.63}{#1}}
\newcommand{\SpecialStringTok}[1]{\textcolor[rgb]{0.73,0.40,0.53}{#1}}
\newcommand{\StringTok}[1]{\textcolor[rgb]{0.25,0.44,0.63}{#1}}
\newcommand{\VariableTok}[1]{\textcolor[rgb]{0.10,0.09,0.49}{#1}}
\newcommand{\VerbatimStringTok}[1]{\textcolor[rgb]{0.25,0.44,0.63}{#1}}
\newcommand{\WarningTok}[1]{\textcolor[rgb]{0.38,0.63,0.69}{\textbf{\textit{#1}}}}
\usepackage{longtable,booktabs,array}
\newcounter{none} % for unnumbered tables
\usepackage{calc} % for calculating minipage widths
% Correct order of tables after \paragraph or \subparagraph
\usepackage{etoolbox}
\makeatletter
\patchcmd\longtable{\par}{\if@noskipsec\mbox{}\fi\par}{}{}
\makeatother
% Allow footnotes in longtable head/foot
\IfFileExists{footnotehyper.sty}{\usepackage{footnotehyper}}{\usepackage{footnote}}
\makesavenoteenv{longtable}
\setlength{\emergencystretch}{3em} % prevent overfull lines
\providecommand{\tightlist}{%
  \setlength{\itemsep}{0pt}\setlength{\parskip}{0pt}}
\usepackage{bookmark}
\IfFileExists{xurl.sty}{\usepackage{xurl}}{} % add URL line breaks if available
\urlstyle{same}
\hypersetup{
  hidelinks,
  pdfcreator={LaTeX via pandoc}}

\author{}
\date{}

\begin{document}

\section{Compte rendu --- Prédiction du risque de défaut (German
Credit)}\label{compte-rendu-pruxe9diction-du-risque-de-duxe9faut-german-credit}

SAMID Oumaima · ENCG Settat · Finance · Partie B --- 9 algorithmes
supervisés

\subsection{1. Sujet et problématique}\label{sujet-et-probluxe9matique}

Une banque doit décider d'accorder ou non un crédit sans connaître à
l'avance la capacité de remboursement de l'emprunteur. Un « mauvais
risque » accordé coûte bien plus qu'un « bon client » refusé : l'erreur
n'est pas symétrique.

Problématique : à partir des caractéristiques d'un dossier (compte,
durée, montant, historique, emploi, âge\ldots), peut-on prédire de façon
fiable qu'un emprunteur sera en défaut, en repérant surtout les mauvais
risques (classe 1, 30 \% des dossiers) ?

Objectif mesurable fixé avant l'accès au test : F1 de la classe «
mauvais risque » ≥ 0,55, à comparer à la référence naïve « tout le monde
est mauvais » (F1 = 0,462 ; exactitude 0,300).

La métrique principale est le F1 de la classe 1, car l'exactitude seule
est trompeuse sur des classes déséquilibrées.

\subsection{2. Jeu de données}\label{jeu-de-donnuxe9es}

Source : German Credit Data (Statlog), dépôt UCI Machine Learning
Repository, 1994 (Prof.~Hofmann, Hambourg). Licence : CC BY 4.0 selon la
page UCI (à vérifier avant remise). Taille : 1 000 dossiers, 20
variables explicatives (7 numériques, 13 catégorielles) plus la cible.
Unité d'observation : un emprunteur / une demande de crédit.

Cible \texttt{risque} : 1 = mauvais risque (300 dossiers, 30 \%), 0 =
bon client (700). Aucune valeur manquante, aucun doublon. La colonne
\texttt{n°} n'est qu'un identifiant : elle est supprimée.

Limites : données anciennes (1994, Allemagne), petit échantillon (200
dossiers de test, dont 60 mauvais), variables parfois codées de façon
grossière, présence de variables sensibles (âge, statut personnel,
nationalité). Adéquation : le jeu convient au sujet (décision d'octroi,
cible binaire déséquilibrée) mais ne garantit pas la validité sur un
portefeuille actuel.

{\def\LTcaptype{none} % do not increment counter
\begin{longtable}[]{@{}
  >{\raggedright\arraybackslash}p{(\linewidth - 4\tabcolsep) * \real{0.3333}}
  >{\raggedright\arraybackslash}p{(\linewidth - 4\tabcolsep) * \real{0.3333}}
  >{\raggedright\arraybackslash}p{(\linewidth - 4\tabcolsep) * \real{0.3333}}@{}}
\toprule\noalign{}
\begin{minipage}[b]{\linewidth}\raggedright
Variable
\end{minipage} & \begin{minipage}[b]{\linewidth}\raggedright
Rôle
\end{minipage} & \begin{minipage}[b]{\linewidth}\raggedright
Remarque
\end{minipage} \\
\midrule\noalign{}
\endhead
\bottomrule\noalign{}
\endlastfoot
\texttt{duree\_mois} & Numérique & Durée du crédit ; 70 valeurs extrêmes
(7,0 \%) \\
\texttt{montant} & Numérique & Montant ; 72 valeurs extrêmes (7,2 \%),
asymétrie à droite \\
\texttt{age} & Numérique & Âge ; 23 valeurs extrêmes (2,3 \%) \\
Variables catégorielles (compte, historique, objet, emploi,
logement\ldots) & Qualitatives & Encodées en one-hot \\
\texttt{risque} & Cible binaire & 1 = mauvais risque (30 \%) \\
\texttt{n°} & Identifiant & Supprimé \\
\end{longtable}
}

\subsection{3. Exploration des
données}\label{exploration-des-donnuxe9es}

Statistiques : 1 000 lignes, aucune valeur manquante ni doublon ;
classes déséquilibrées (70 \% / 30 \%). Les montants sont fortement
asymétriques à droite.

\begin{itemize}
\tightlist
\item
  Graphique 1 --- répartition de la cible : 30 \% de mauvais risques
  seulement ; un modèle qui répond toujours « bon client » aurait 70 \%
  d'exactitude sans rien détecter, d'où le choix du F1 et d'un découpage
  stratifié.
\item
  Graphique 2 --- distribution du montant et de la durée : queue longue
  à droite, ce qui justifie la variable \texttt{log\_montant} et la
  normalisation pour les modèles sensibles à l'échelle.
\item
  Graphique 3 --- variables explicatives selon la cible (durée, montant,
  âge, compte) : les crédits plus longs et plus élevés sont plus souvent
  des mauvais risques ; l'effet est réel mais partiel, d'où des classes
  qui se chevauchent.
\end{itemize}

Qualité des données : valeurs extrêmes détectées par la règle IQR × 1,5
(durée 70, montant 72, âge 23) ; ce sont des dossiers plausibles, ils
sont conservés.

\subsection{4. Préparation des
données}\label{pruxe9paration-des-donnuxe9es}

Traitements : suppression de \texttt{n°} (identifiant) et séparation de
la cible \texttt{risque} avant toute transformation, pour éviter toute
fuite. Les variables catégorielles sont encodées en one-hot (63
variables au total). Les valeurs extrêmes sont conservées.

Variables créées :
\texttt{mensualite\_estimee\ =\ montant\ /\ duree\_mois} (poids de
l'échéance) et \texttt{log\_montant\ =\ log1p(montant)} (réduit
l'asymétrie).

Prévention des fuites : le découpage train/test est fait en premier ; le
\texttt{StandardScaler} est placé dans un Pipeline, donc ajusté sur les
seuls plis d'entraînement, et seulement pour la régression logistique,
le KNN et le SVM. Les modèles à base d'arbres n'en ont pas besoin. Le
seuil de décision est choisi par validation croisée, jamais sur le test.

\subsection{5. Algorithme principal : régression
logistique}\label{algorithme-principal-ruxe9gression-logistique}

Principe : modèle linéaire qui estime la probabilité de défaut ; chaque
variable a un coefficient lisible, ce qui compte en crédit où la
décision doit être expliquée. Formule :

\begin{Shaded}
\begin{Highlighting}[]
\NormalTok{P(y=1 }\FunctionTok{\textbackslash{}mid}\NormalTok{ x) = }\FunctionTok{\textbackslash{}frac}\NormalTok{\{1\}\{1 + e\^{}\{{-}(}\FunctionTok{\textbackslash{}beta}\NormalTok{\_0 + }\FunctionTok{\textbackslash{}beta}\NormalTok{\^{}\{}\FunctionTok{\textbackslash{}top}\NormalTok{\} x)\}\}}
\end{Highlighting}
\end{Shaded}

Les coefficients minimisent la perte logistique pénalisée L2 :
\texttt{C} est l'inverse de la force de pénalisation.

Hyperparamètres retenus : \texttt{C\ =\ 0,2} (régularisation assez forte
pour 63 variables et 800 observations), \texttt{max\_iter\ =\ 1500},
variables standardisées.

Raisons du choix : meilleur F1 test des 9 modèles (0,625), AUC élevée
(0,810), interprétable, stable et rapide. Les huit autres algorithmes
(KNN, Naive Bayes, arbre, forêt, SVM RBF, Gradient Boosting, AdaBoost,
XGBoost) servent de comparaison.

\subsection{6. Protocole expérimental}\label{protocole-expuxe9rimental}

Découpage stratifié 80 / 20 avec \texttt{random\_state\ =\ 42} : 800
dossiers d'entraînement, 200 de test (60 mauvais risques). Référence :
classifieur naïf « tout mauvais » (F1 = 0,462).

Validation : 5 plis stratifiés sur le train uniquement, utilisés pour
comparer les modèles, régler le seuil de décision (probabilités
hors-pli) et balayer \texttt{C} et \texttt{k}. Le test n'est utilisé
qu'une fois, pour le bilan final.

Métriques : F1 de la classe 1 (principale), rappel, précision,
exactitude, AUC ; matrice de confusion pour l'analyse des erreurs.

\end{document}
